% block: computational-physics-heading page:63 type:section
\subsection{{Computational Physics}}\label{sec:computational-physics}
% block: electronic-structure-heading page:63 type:subsection
\subsubsection{{Large-Scale Electronic Structure Calculation}}\label{subsec:electronic-structure}
{Molecular dynamics simulations in large scale systems can be efficiently computed by using the one-body density matrix $\rho$, which is based on the fact that any physical quantity $\langle X\rangle$ can be obtained by}
\[\langle X\rangle=\operatorname{Tr}(\rho X)=\iint\rho(\boldsymbol{r},\boldsymbol{r}')X(\boldsymbol{r}',\boldsymbol{r})\,\mathrm d\boldsymbol{r}\mathrm d\boldsymbol{r}'\]
{and the corresponding discretization form is written as}
\[\sum_{i,j}\rho_{ij}X_{ji}.\]
{It is important to note here that if $X$ is a short-range operator, we only need the short-range behavior of the density matrix. As an extreme example, if $X$ is a diagonal matrix (i.e., off-diagonal elements are zeros), then we only need the diagonal element of the density matrix because $\sum_{i,j}\rho_{ij}X_{ji}=\sum_i\rho_{ii}X_{ii}$, leading to very efficient computation.}
% block: density-matrix-green-function pages:63,64 type:ordinary-paragraph
{In what follows, we explain the density matrix $\rho$ in detail. It is known that the $(i,j)$ element of the density matrix $\rho$ is given by}
\begin{equation}\label{eq:density-matrix-green-integral}
\rho_{ij}=-\frac1\pi\int_{-\infty}^{\infty}\operatorname{Im}G_{ij}(\epsilon)f\left(\frac{\epsilon-\mu}{K_BT}\right)\,\mathrm d\epsilon,
\end{equation}
{where $f(x)$ is the Fermi distribution function, $\mu$ is a chemical potential that may be determined so that the total number of electrons equals the sum of the diagonal elements of the density matrix. Here, $G_{ij}(\epsilon)$ is the $(i,j)$ element of the following matrix:}
\begin{equation}\label{eq:green-function-resolvent}
G(\epsilon)=[(\epsilon+i\delta)I-H]^{-1},
\end{equation}
{where $i$ is the unit imaginary number, $\delta$ is an (infinitely) small positive number ($\delta\to0_+$), and $I$ and $H$ are the identity matrix and a Hamiltonian matrix. The matrix $G$ is referred to as Green's function. The $(i,j)$ element of the Hamiltonian matrix $H$ may be given by the following real space integration:}
\[H_{ij}=\int f_i(\boldsymbol{r})\left(\frac{-\hbar}{2m}\Delta+V(\boldsymbol{r})\right)f_j(\boldsymbol{r})\,\mathrm d\boldsymbol{r},\]
{where $m$ and $\hbar$ are the mass of the electron and the (reduced) Planck constant, $\Delta$ is the Laplacian. $f_j(\boldsymbol{r})$'s are prescribed functions that are linearly independent. $V(\boldsymbol{r})$ is a prescribed function regarding the interactions between electrons and nuclei. For the details, see, e.g., \BookCite{184}, \BookCite{190} and the references therein.}
% block: electronic-shifted-systems page:64 type:ordinary-paragraph
{We now describe how linear systems arise from the above computation. Let $\boldsymbol{e}_i$ be the $i$th unit vector. Then from \eqref{eq:green-function-resolvent}, the $(i,j)$ element of $G(\epsilon)$ is calculated by}
\[G_{ij}(\epsilon)=\boldsymbol{e}_i^{\top}[(\epsilon+i\delta)I-H]^{-1}\boldsymbol{e}_j.\]
{Therefore, letting $\boldsymbol{x}_j=[(\epsilon+i\delta)I-H]^{-1}\boldsymbol{e}_j$ yields}
\[[(\epsilon+i\delta)I-H]\boldsymbol{x}_j=\boldsymbol{e}_j,\]
{which is a linear system. In computational physics, the Hamiltonian matrix $H$ is Hermitian or real symmetric. If $H$ is real symmetric, the coefficient matrix $(\epsilon+i\delta)I-H$ is complex symmetric. Furthermore, for computing the integration in \eqref{eq:density-matrix-green-integral}, a suitable numerical quadrature formula is used, which requires $G_{ij}(\epsilon_k)$ for $k=1,2,\ldots,m$. Thus, for computing the $(i,j)$ element of the density matrix $\rho_{i,j}$, we need to solve linear systems of the form}
\[(A+\sigma_kI)\boldsymbol{x}_j^{(k)}=\boldsymbol{b},\quad k=1,2,\ldots,m,\]
{where $A:=i\delta I-H$, $\boldsymbol{b}:=\boldsymbol{e}_j$ and $\sigma_k:=\epsilon_k$. These are known as shifted linear systems, which are efficiently solved by Krylov subspace methods by utilizing the shift-invariance of Krylov subspace. Krylov subspace methods for shifted linear systems are described in Chap.\ \ref{chap:shifted-linear-systems}.}
% block: lattice-qcd-heading pages:64,65 type:subsection
\subsubsection{{Lattice Quantum Chromodynamics}}\label{subsec:lattice-qcd}
{Quantum chromodynamics (QCD) is the theory of the strong interaction between quarks and gluons, and lattice QCD is a non-perturbative approach to solving QCD, where the theory is formulated on a lattice in space and time. It is known that the most time-consuming part of lattice QCD is solving linear systems}
\[A\boldsymbol{x}=\boldsymbol{b},\]
{which are often written as $A\psi=\phi$ and the solution is used for obtaining quark propagators. For the details of the quark propagators including a summary of lattice QCD, see, e.g., \BookCite{69}. In the following, we describe the structure of the Wilson fermion matrix $A$.}
% block: wilson-hopping-matrix page:65 type:ordinary-paragraph
{The Wilson fermion matrix is written as $A=I-\kappa D$, where $I$ is the identity matrix, and $\kappa$ is a real nonnegative parameter, referred to as the hopping parameter. Matrix $D$ is called the hopping matrix. We now explain the structure of the hopping matrix.}
{Let us consider a four-dimensional hypercubic lattice which can be regarded as an equispaced grid of the hypercube. Each grid point can be written as a vector $\boldsymbol{x}$ such that}
\[\boldsymbol{x}\in\Omega:=\{(x_1,x_2,x_3,x_4):x_1,x_2,x_3,x_4\in\{1,2,\ldots,N\}\}.\]
{For simplicity, we assume that $x_1,x_2,x_3,x_4$ run from 1 to the same number $N$. Then it is easy to see that the number of elements of $\Omega$ is $N^4$. Using $\boldsymbol{x},\boldsymbol{y}\in\Omega$, the hopping matrix $D$ is written by block matrix form}
\begin{equation}\label{eq:wilson-hopping-block}
D(\boldsymbol{x};\boldsymbol{y})=\sum_{\mu=1}^4\left\{[(I_4-\gamma_\mu)\otimes U_\mu(\boldsymbol{x})]\delta_{\boldsymbol{x},\boldsymbol{y}-\boldsymbol{e}_\mu}
+[(I_4+\gamma_\mu)\otimes U_\mu^{\mathrm H}(\boldsymbol{x}-\boldsymbol{e}_\mu)]\delta_{\boldsymbol{x},\boldsymbol{y}+\boldsymbol{e}_\mu}\right\}.
\end{equation}
% block: wilson-full-block-matrix page:65 type:ordinary-paragraph
{Here, $\delta_{\boldsymbol{x},\boldsymbol{y}}$ is the Kronecker delta, i.e., $\delta_{\boldsymbol{x},\boldsymbol{y}}=1$ if $\boldsymbol{x}=\boldsymbol{y}$ and $\delta_{\boldsymbol{x},\boldsymbol{y}}=0$ if $\boldsymbol{x}\ne\boldsymbol{y}$. $D(\boldsymbol{x};\boldsymbol{y})$ is a small matrix of the size $12\times12$ and the hopping matrix $D$ is given as follows:}
\[\begin{aligned}
D&=\begin{bmatrix}D_{1,1}&D_{1,2}&\cdots&D_{1,N^4}\\D_{2,1}&D_{2,2}&\cdots&D_{2,N^4}\\\vdots&\vdots&\ddots&\vdots\\D_{N^4,1}&D_{N^4,2}&\cdots&D_{N^4,N^4}\end{bmatrix}\\
&=\begin{bmatrix}
D(1,1,1,1;1,1,1,1)&D(1,1,1,1;1,1,1,2)&\cdots&D(1,1,1,1;N,N,N,N)\\
D(1,1,1,2;1,1,1,1)&D(1,1,1,2;1,1,1,2)&\cdots&D(1,1,1,2;N,N,N,N)\\
\vdots&\vdots&\ddots&\vdots\\
D(N,N,N,N;1,1,1,1)&D(N,N,N,N;1,1,1,2)&\cdots&D(N,N,N,N;N,N,N,N)
\end{bmatrix}.
\end{aligned}\]
{The size of the hopping matrix $D$ is thus $12N^4\times12N^4$. If $N=64$, then the number of unknowns for the linear systems is about 200 million! The size grows much more rapidly than that arising from the discretization of three-dimensional partial differential equations in Sect.\ \ref{subsec:nonsymmetric-difference}.}
% block: wilson-gamma-and-pauli page:66 type:ordinary-paragraph
{The details of \eqref{eq:wilson-hopping-block} are explained as follows: $I_4$ is the 4-by-4 identity matrix, $\delta_{\boldsymbol{x},\boldsymbol{y}}$ is a function ($\delta:\mathbb{R}^4\times\mathbb{R}^4\to\mathbb{R}$), where $\delta_{\boldsymbol{x},\boldsymbol{y}}=0$ if $\boldsymbol{x}\ne\boldsymbol{y}$, and $\delta_{\boldsymbol{x},\boldsymbol{y}}=1$ if $\boldsymbol{x}=\boldsymbol{y}$. The symbol $\boldsymbol{e}_\mu$ is the transpose of the $\mu$th unit vector, or equivalently the $\mu$th row of the $4\times4$ identity matrix $I_4$. For each $\mu$, the symbol $\gamma_\mu$ is a matrix given by}
\[\gamma_1=\sigma_y\otimes\sigma_x,\quad\gamma_2=\sigma_y\otimes\sigma_y,\quad\gamma_3=\sigma_y\otimes\sigma_z,\quad\gamma_4=\sigma_z\otimes I_2,\]
{where $I_2$ is the 2-by-2 identity matrix, and $\sigma_x,\sigma_y$, and $\sigma_z$ are the Pauli matrices, i.e.,}
\[\sigma_x=\begin{bmatrix}0&1\\1&0\end{bmatrix},\quad\sigma_y=\begin{bmatrix}0&-i\\i&0\end{bmatrix},\quad\sigma_z=\begin{bmatrix}1&0\\0&-1\end{bmatrix}.\]
% block: wilson-gauge-field page:66 type:ordinary-paragraph
{Finally, for each $\mu$ and $\boldsymbol{x}$, the symbol $U_\mu(\boldsymbol{x})(=U_\mu([x_1,x_2,x_3,x_4])$ for $\mu=1,2,3,4$ belongs to $\operatorname{SU}(3)$, i.e., a 3-by-3 unitary matrix whose determinant is one. $U_\mu(\boldsymbol{x})$ corresponds to the background gauge field. $U_\mu(\boldsymbol{x})$ is called ``cold'' if $U_\mu(\boldsymbol{x})=I_3$ and ``hot'' if the $U_\mu(\boldsymbol{x})\in\operatorname{SU}(3)$ is randomly chosen. Below are some examples for $D(\boldsymbol{x},\boldsymbol{y})$ when $N=2$:}
\[\begin{aligned}
D_{1,1}&=D(1,1,1,1;1,1,1,1)=O,\\
D_{1,2}&=D(1,1,1,1;1,1,1,2)=(I_4-\gamma_4)\otimes U_4([1,1,1,1]),\\
D_{1,3}&=D(1,1,1,1;1,1,2,1)=(I_4-\gamma_3)\otimes U_3([1,1,1,1]),\\
D_{1,4}&=D(1,1,1,1;1,1,2,2)=O,\\
D_{1,5}&=D(1,1,1,1;1,2,1,1)=(I_4-\gamma_2)\otimes U_2([1,1,1,1]),\\
D_{1,6}&=D(1,1,1,1;1,2,1,2)=O,\\
D_{1,7}&=D(1,1,1,1;1,2,2,1)=O,\\
D_{1,8}&=D(1,1,1,1;1,2,2,2)=O,\\
D_{1,9}&=D(1,1,1,1;2,1,1,1)=(I_4-\gamma_1)\otimes U_1([1,1,1,1]),\\
&\vdots\\
D_{2,1}&=D(1,1,1,2;1,1,1,1)=(I_4+\gamma_4)\otimes U_4^{\mathrm H}([1,1,1,1]),\\
D_{2,2}&=D(1,1,1,2;1,1,1,2)=O,\\
D_{2,3}&=D(1,1,1,2;1,1,2,1)=O,\\
D_{2,4}&=D(1,1,1,2;1,1,2,2)=(I_4-\gamma_3)\otimes U_3([1,1,1,2]),\\
&\vdots\\
D_{16,16}&=D([2,2,2,2;2,2,2,2])=O.
\end{aligned}\]
{Note that $D_{2,1}=(I_4+\gamma_4)\otimes U_4^{\mathrm H}([1,1,1,2]-\boldsymbol{e}_4)=(I_4+\gamma_4)\otimes U_4^{\mathrm H}([1,1,1,1])$. Matrix $D$ is neither symmetric nor Hermitian, but it has a hidden symmetry as described next. Let}
% block: gamma-five-symmetry page:67 type:ordinary-paragraph
\[\gamma_5=\sigma_x\otimes I_2.\]
{Then, the hopping matrix $D$ has the following property:}
\begin{equation}\label{eq:gamma-five-hermiticity}
\Gamma_5D=D^{\mathrm H}\Gamma_5,
\end{equation}
{where $\Gamma_5=I_{N^4}\otimes\gamma_5\otimes I_3$. This property is known as $\gamma_5$-symmetric (or $\gamma_5$-Hermitian). From Theorem\ \ref{thm:kronecker-product-properties}-(5), matrix $\gamma_5$ is real symmetric and matrix $\Gamma_5$ is also real symmetric, and thus we have $\Gamma_5=\Gamma_5^{\mathrm H}$. Then, we can rewrite \eqref{eq:gamma-five-hermiticity} as}
\[\Gamma_5D=(\Gamma_5D)^{\mathrm H}.\]
{This means that $\Gamma_5D$ is Hermitian. From Corollary\ \ref{cor:kronecker-inverse}, we have $\Gamma_5^{-1}=I_{N^4}^{-1}\otimes\gamma_5^{-1}\otimes I_3^{-1}=I_{N^4}\otimes\gamma_5\otimes I_3=\Gamma_5$. Thus we have another representation of \eqref{eq:gamma-five-hermiticity}:}
\begin{equation}\label{eq:gamma-five-similarity}
\Gamma_5D\Gamma_5=D^{\mathrm H}.
\end{equation}
% block: gamma-five-block-verification page:67 type:ordinary-paragraph
{Now, let us confirm \eqref{eq:gamma-five-similarity} for $N=2$. Let $P=\gamma_5\otimes I_3$. Then the $(i,j)$ block element of $D$ is}
\[\begin{aligned}
(\Gamma_5D\Gamma_5)_{i,j}
&=\left(\begin{bmatrix}P&&\\&\ddots&\\&&P\end{bmatrix}
\begin{bmatrix}D_{1,1}&\cdots&D_{1,16}\\\vdots&\ddots&\vdots\\D_{16,1}&\cdots&D_{16,16}\end{bmatrix}
\begin{bmatrix}P&&\\&\ddots&\\&&P\end{bmatrix}\right)_{i,j}\\
&=\left(\begin{bmatrix}PD_{1,1}P&\cdots&PD_{1,16}P\\\vdots&\ddots&\vdots\\PD_{16,1}P&\cdots&PD_{16,16}P\end{bmatrix}\right)_{i,j}\\
&=PD_{i,j}P.
\end{aligned}\]
{As an example, consider $(i,j)=(1,2)$. Then, from Theorem\ \ref{thm:kronecker-product-properties}-(5) and (6), it follows that}
\[\begin{aligned}
(\Gamma_5D\Gamma_5)_{1,2}&=PD_{1,2}P\\
&=(\gamma_5\otimes I_3)[(I_4-\gamma_4)\otimes U_4([1,1,1,1])](\gamma_5\otimes I_3)\\
&=\gamma_5(I_4-\gamma_4)\gamma_5\otimes U_4([1,1,1,1])\\
&=(I_4-\gamma_5\gamma_4\gamma_5)\otimes U_4([1,1,1,1])\\
&=(I_4+\gamma_4)\otimes U_4([1,1,1,1])\\
&=[(I_4+\gamma_4)^{\mathrm H}\otimes U_4^{\mathrm H}([1,1,1,1])]^{\mathrm H}\\
&=[(I_4+\gamma_4)\otimes U_4^{\mathrm H}([1,1,1,1])]^{\mathrm H}\\
&=D_{2,1}^{\mathrm H}.
\end{aligned}\]
% block: lattice-eigenvalue-problems page:68 type:ordinary-paragraph
{Similarly, we can confirm that $(\Gamma_5D\Gamma_5)_{i,j}=D_{j,i}^{\mathrm H}$ for the other $i,j$'s. Thus we see that \eqref{eq:gamma-five-similarity} holds true for the case $N=2$. In lattice QCD, there is a need to compute some small eigenvalues in magnitude. In this case, the shift-and-invert Lanczos (Arnoldi) method may be used. The method requires solving linear systems at each iteration. See \BookCite{12} for algorithms to compute eigenvalue problems. A thick-restart Lanczos type method for eigenvalue problems arising in lattice QCD was proposed in \BookCite{108}.}
